\documentclass[11pt]{article}

%==============================================================================
%  LPS / PS-LPS Program: Plan of Action and Phase Sequencing
%  A dependency-ordered roadmap across validation, synchronization, dimension
%  regularization, covariance + simplicial complex, and occupancy decomposition.
%  All figures original. References point to the program's own design notes.
%==============================================================================

\usepackage[margin=1in]{geometry}
\usepackage{amsmath,amssymb}
\usepackage{graphicx}
\usepackage{xcolor}
\usepackage{booktabs}
\usepackage{enumitem}
\usepackage{microtype}
\microtypesetup{expansion=false}   % avoid pdfTeX font-expansion error in this build
\usepackage{caption}
\usepackage[most]{tcolorbox}
\usepackage{datetime2}
\DTMsetup{showzone=false}   % timestamp in US Eastern time (build sets TZ=America/New_York)
\usepackage[hidelinks]{hyperref}
\usepackage[section]{placeins}
\widowpenalty=10000
\clubpenalty=10000
\raggedbottom

\definecolor{ink}{HTML}{24303D}
\definecolor{accent}{HTML}{2563EB}
\definecolor{accent3}{HTML}{15803D}
\captionsetup{font=small,labelfont=bf}
\setlist{topsep=3pt,itemsep=2pt,parsep=1pt}

\newtcolorbox{gate}{
  enhanced, colback=accent3!5, colframe=accent3!55!black, boxrule=0.4pt, arc=2pt,
  left=7pt, right=7pt, top=4pt, bottom=4pt, title=Go/no-go gate,
  fonttitle=\bfseries\sffamily\small, coltitle=accent3!45!black,
  attach title to upper={\quad}}
\newtcolorbox{note}{
  enhanced, colback=accent!4, colframe=accent!55!black, boxrule=0.4pt, arc=2pt,
  left=7pt, right=7pt, top=4pt, bottom=4pt, title=Why this ordering,
  fonttitle=\bfseries\sffamily\small, coltitle=accent!50!black,
  attach title to upper={\quad}}

\newcommand{\field}[1]{\smallskip\noindent\textbf{\textsf{#1.}}\ }
\newcommand{\fig}[3]{%
  \begin{figure}[htbp]\centering
  \includegraphics[width=#2]{figures_ldr/#1}%
  \caption{#3}\label{fig:#1}\end{figure}}

\title{\textbf{LPS / PS-LPS Program: Plan of Action and Phase Sequencing}\\[4pt]
\large A dependency-ordered roadmap across validation, synchronization,
local-dimension regularization, local covariance and the simplicial-complex
substrate, synthetic 16S data, and the occupancy / quasi-equilibrium decomposition}
\author{\small Internal program plan --- LPS / \texttt{geosmooth} program\\[2pt]
\footnotesize source: \nolinkurl{/Users/pgajer/current_projects/geosmooth/project_briefs/}\\[1pt]
\footnotesize\nolinkurl{lps_program_plan_of_action_2026-06-10.tex}}
\date{\DTMnow\ \textnormal{\footnotesize ET}}

\begin{document}
\maketitle

\begin{abstract}\noindent
The audit of the LPS/PS-LPS effort opened into a layered research program. This
note fixes the \emph{order} in which its pieces should be built, with the
dependencies made explicit so the sequence is a consequence of what each piece
needs rather than a preference. The spine follows the components as listed ---
validation first, then the estimator extensions, then the compositional/covariance
layer and its simplicial-complex substrate, and finally the occupancy /
quasi-equilibrium decomposition --- with three structural refinements: (i) the
synthetic-data and VALENCIA-exploration workstream is pulled \emph{forward} as a
cross-cutting enabler, because it is a dependency of almost every later test, not a
last-bundle item; (ii) PS-LPS synchronization and local-dimension regularization run
\emph{in parallel} as independent extensions of the validated base, then combine;
and (iii) the simplicial complex (the nerve of the chart cover) is named as the
\emph{integration} phase that unifies synchronization, dimension, and covariance,
not as a sub-part of any one of them. Each phase is specified by its objective,
inputs (dependencies), key deliverables (the existing design notes are its specs),
and a go/no-go gate; the whole program runs under the two-agent audit-independence
workflow. The immediate next actions are the two audits already queued: the LPS
Tier-0 run and the LPS binary GM/FF run.
\par\medskip\noindent\footnotesize\textit{On the figures.} Both figures are
original schematics drawn for this note.
\end{abstract}

{\small\tableofcontents}

%==============================================================================
\section{How to read this, and the one rule that wraps every phase}\label{sec:read}
%==============================================================================

This is a sequencing document, not a methods document: the methods live in the
design notes referenced at the end, and each phase below points to the note that
specifies it. The job here is only to say what is done first, second, and so on,
and why the order is forced.

One rule wraps every phase. Each phase is executed under the \emph{two-agent
audit-independence workflow}: an implementer builds, an independent auditor
validates against a fixed charter (the auditor's scope is not set by the
implementer, the handoff carries evidence and admitted limitations but not the
audit's questions, and ``report rendered'' is never ``phase accepted''). A phase is
not complete --- and the next phase does not start --- until its go/no-go gate is
passed by that independent audit or adjudicated by the investigator. This is what
keeps a fast-moving, multi-layer program from accumulating unvalidated machinery.

\begin{note}
The ordering is driven by \emph{dependencies}, not by importance or interest. A
phase may begin only once every input it draws an arrow from in
Figure~\ref{fig:fig_program_dag.pdf} has passed its gate. Two phases with no
dependency between them run in parallel. The capstone is last because it depends on
the most other pieces, not because it is an afterthought.
\end{note}

%==============================================================================
\section{The components, and how the dependencies fall}\label{sec:components}
%==============================================================================

The program has five technical components plus one cross-cutting enabler and one
terminal goal. Stated as the dependency structure of
Figure~\ref{fig:fig_program_dag.pdf}:

\begin{itemize}[leftmargin=1.4em]
\item \textbf{Phase 0 --- LPS validation.} The base smoother, proven correct
(polynomial reproduction, the linear-smoother identity, df, honest cross-validation,
boundary behaviour, binary calibration, degenerate-geometry robustness). It
\emph{emits} a toolkit --- the smoother matrix $S$, $\mathrm{tr}\,S$, analytic
leave-one-out / GCV, and uncertainty bands --- that later phases reuse. Foundation
for everything.
\item \textbf{Cross-cutting --- the data enabler.} Synthetic non-manifold
generators with known ground-truth dimension/components, maturing from toy
stratified objects to VALENCIA-realistic 16S generators via the E1/E2/E3
exploration. A \emph{dependency} of the dimension field, the covariance layer, and
the occupancy capstone --- hence front-loaded, not last.
\item \textbf{Phase 1a --- PS-LPS synchronization.} Prediction-overlap coupling of
local models. An independent extension of the validated base; reuses the Phase-0
toolkit for the df-matched comparison and exact $\lambda_{\mathrm{sync}}$ selection.
\item \textbf{Phase 1b --- local-dimension regularization.} The $\ell_1$/TV
dimension field with exact graph-cut solve. Independent of 1a; needs the validated
base and the synthetic generators.
\item \textbf{Phase 2a --- local covariance / association.} The chart-aware
covariance/precision layer on the compositional ($L^\infty$) geometry; needs the
base and the realistic 16S generators.
\item \textbf{Phase 2b --- the simplicial complex (nerve): integration.} The nerve
of the chart cover, unifying 1a + 1b + 2a into one consistency objective over the
overlaps; depends on all three.
\item \textbf{Phase 3 (capstone) --- occupancy / quasi-equilibrium / support
decomposition.} Finding pure components and transition phases; builds on the
covariance and the nerve and the realistic generators.
\item \textbf{Terminal --- the real 16S / omics application.} The actual estimation
on real data, which the whole program exists to serve.
\end{itemize}

\fig{fig_program_dag.pdf}{\linewidth}{The program dependency DAG. Arrows point from
prerequisite to dependent: solid $=$ direct dependency, dashed $=$ reuses the
Phase-0 toolkit, dotted $=$ uses the synthetic data. PS-LPS (1a) and the dimension
field (1b) are parallel; the nerve (2b) integrates 1a+1b+2a; the occupancy capstone
(Phase 3) and the terminal application sit downstream of the integration.}

\fig{fig_program_timeline.pdf}{\linewidth}{The same plan as parallel tracks over
time. Validation runs first; the data enabler runs throughout; the two Phase-1
extensions run in parallel on the validated base; the covariance and nerve-integration
tracks follow; the occupancy capstone and then the real-data application close the
program. Diamonds are independent-audit go/no-go gates between phases.}

%==============================================================================
\section{Phase 0 --- LPS validation (the foundation that emits the tools)}\label{sec:phase0}
%==============================================================================

\field{Objective} Prove the base smoother is what it claims, and produce the
reusable estimator toolkit.
\field{Inputs} None (the foundation).
\field{Deliverables} The Tier-0 correctness battery (polynomial reproduction; the
linear-smoother identity $\hat y = Sy$ with $\mathrm{tr}\,S$ and analytic LOO/GCV;
boundary behaviour; consistency; binary recovery and calibration; CV no-leakage;
degenerate geometry), then the rest of the experimental plan (honest nested/grouped
CV; bandwidth/kernel semantics; binary-mode hygiene; uncertainty bands). Spec: the
\emph{experimental plan} note.
\field{Gate} \begin{gate} The Tier-0 battery passes deterministically; the
linear-smoother toolkit ($S$, df, LOO/GCV, bands) is implemented and tested. No
later phase begins until this gate is green --- it is both the correctness gate and
the source of the tools the df-matched, GCV, and band computations downstream
require.\end{gate}
\field{Unblocks} Everything --- and specifically hands PS-LPS its df-matched
comparison and exact $\lambda_{\mathrm{sync}}$ selection.

\field{Immediate next action} The two audits already queued are the first concrete
execution of this phase: (1) the \textbf{LPS Tier-0 run} (does the correctness
battery actually pass, and is the toolkit correct?), and (2) the \textbf{LPS binary
GM/FF run} (the binary-mode behaviour, with the fallback-telemetry and
selection-metric fixes from the earlier results audit). Both are Phase-0 evidence;
neither downstream phase should start before they are adjudicated.

%==============================================================================
\section{Cross-cutting --- the synthetic-data and VALENCIA-exploration enabler}\label{sec:data}
%==============================================================================

\field{Objective} Supply known-ground-truth, increasingly realistic non-manifold
datasets, because every downstream test needs an answer key the real data cannot
give.
\field{Inputs} Phase 0 (charts), and --- for the realistic versions --- the VALENCIA
relative-abundance matrix.
\field{Deliverables} Early: the toy stratified generators (the implementation
notes' DGP library) needed by Phase 1b. Maturing: the VALENCIA E1 (enumerate naive
strata), E2 (mass concentration near lower faces), E3 (parameter bundle) exploration,
producing the realistic stratum-mixture generators with known components that feed
Phases 2 and 3. Spec: the \emph{implementation notes}.
\field{Gate} \begin{gate} Toy generators ready before Phase 1b's boundary/scale
tests. Realistic generators pass a posterior-predictive fidelity check (synthetic
marginals reproduce the real VALENCIA marginals) before any Phase-2/3 test built on
them is trusted.\end{gate}
\field{Unblocks} Phase 1b (toy), Phase 2a and Phase 3 (realistic).

\begin{note}
This is the one item promoted out of the original ordering. In the initial list,
synthetic models sat bundled with the covariance layer near the end. But the
dimension-field tests (1b), the covariance tests (2a), and the occupancy capstone
(3) all consume known-truth synthetic data, and the realistic versions need the
VALENCIA exploration to exist first. A dependency cannot trail the things that
depend on it, so the data workstream leads and runs throughout.
\end{note}

%==============================================================================
\section{Phase 1 --- two parallel extensions of the validated base}\label{sec:phase1}
%==============================================================================

PS-LPS (1a) and the local-dimension field (1b) are \emph{independent} of each
other: neither uses the other's output. Both depend only on Phase 0 and the
synthetic generators. They therefore run in parallel and are combined later, in
Phase 2b --- which is exactly the ``test independently, then combine'' structure the
program already anticipated.

\subsection{Phase 1a --- PS-LPS prediction synchronization}
\field{Objective} Validate, and put on firm footing, the synchronization of
overlapping local models through prediction-overlap regularization.
\field{Inputs} Phase 0 (the validated base \emph{and} the toolkit).
\field{Deliverables} Exploit the fact that PS-LPS is linear in the response:
implement exact effective df $=\mathrm{tr}\,S(\lambda_{\mathrm{sync}})$ and analytic
LOO/GCV in place of the $K$-fold materialised CV, and compare PS-LPS to LPS at
\emph{matched df} (not matched $\lambda$); fix the ridge-path discontinuity at
$\lambda_{\mathrm{sync}}=0$; make the synchronization graph (currently a fixed
kNN-rank threshold) and its weights explicit; characterize the estimand and the
$\lambda_{\mathrm{sync}}\to\infty$ limit; add uncertainty bands. Spec: the PS-LPS
sections of the audit memo and the chart-aware synchronization.
\field{Gate} \begin{gate} On synthetic data, PS-LPS improves over a \emph{df-matched}
LPS where the geometry predicts it should; $\lambda_{\mathrm{sync}}$ selection is
well-founded via $\mathrm{tr}\,S$ (no longer railing to a grid boundary for lack of
a df-aware criterion); the $\lambda_{\mathrm{sync}}=0$ gate reproduces LPS exactly
and the path is continuous.\end{gate}
\field{Unblocks} Phase 2b (integration over the nerve).

\subsection{Phase 1b --- local-dimension regularization}
\field{Objective} Validate the $\ell_1$/total-variation local-dimension field and
its exact graph-cut solver.
\field{Inputs} Phase 0; the toy stratified generators.
\field{Deliverables} The test battery T1--T9: graph-cut exactness, $\ell_1$-keeps-
the-boundary vs $\ell_2$-blurs, denoising, the boundary safeguard, scale/persistence,
the $\xi$-vs-$u$ coordinate trap, the support cap and soft face, the frontier
(no impossible islands), and downstream Truth-RMSE. Spec: the implementation notes.
\field{Gate} \begin{gate} The implementation notes' single go/no-go: \emph{(the E2
nominal$-$effective gap is large)} $\wedge$ \emph{(T9 improves Truth-RMSE on the
stratum-mixture generator)}. If either fails, the within-face dimension field is not
carried to production; the cheaper soft-face-and-cap refinement is kept
instead.\end{gate}
\field{Unblocks} Phase 2b.

%==============================================================================
\section{Phase 2 --- the compositional covariance, and its integration}\label{sec:phase2}
%==============================================================================

\subsection{Phase 2a --- local covariance / association on the simplex}
\field{Objective} Build and validate the local covariance/precision (the
``shape'' of the local structure) on the boundary-safe compositional geometry.
\field{Inputs} Phase 0; the realistic 16S generators; a compositional front end
(a boundary-aware log-ratio / Aitchison chart with structural-zero and soft-face
handling --- a prerequisite this phase must establish first).
\field{Deliverables} Recovery of a known local covariance on synthetic 16S; the
two-layer association object (co-occurrence support layer $+$ conditional-abundance
precision layer) with the soft-face and $\ell_1$-fusion refinements. Spec: the
chart-aware association design and its addendum.
\field{Gate} \begin{gate} The covariance/precision is recovered on synthetic data
with known structure; the boundary-safe path never reintroduces a pseudocount; the
$\ell_1$ fusion preserves a planted reorganization across a community-state boundary
that $\ell_2$ would smooth away.\end{gate}
\field{Unblocks} Phase 2b and the capstone.

\subsection{Phase 2b --- the simplicial complex (nerve): integration}
\field{Objective} Unify the three preceding layers --- prediction synchronization
(1a), the dimension/size field (1b), and the covariance (2a) --- on the nerve of the
chart cover, as one consistency objective over the overlaps.
\field{Inputs} Phases 1a, 1b, 2a.
\field{Deliverables} The nerve-substrate program, minimal-first: replace PS-LPS's
pairwise overlap graph with the nerve $1$-skeleton of variable-size charts; add
$2$-simplices (higher-order synchronization at junctions); co-regularize chart size
and dimension over the nerve; and, if it pays, the single cochain objective
combining prediction, covariance, and dimension. Spec: the chart-nerve substrate
note.
\field{Gate} \begin{gate} On synthetic data, the nerve $1$-skeleton of variable-size
charts beats the fixed kNN overlap graph on held-out accuracy; higher-order
($2$-simplex) synchronization either measurably helps (at junctions) or is dropped;
the integrated objective is validated on known truth.\end{gate}
\field{Unblocks} Phase 3 --- it supplies the basin-and-corridor skeleton the
capstone runs on.

%==============================================================================
\section{Phase 3 (capstone) --- occupancy / quasi-equilibrium / support}\label{sec:phase3}
%==============================================================================

\field{Objective} Find the pure components (recurrent ecological basins) and the
transition phases between them.
\field{Inputs} Phase 2a (covariance), Phase 2b (the nerve / basin-corridor
skeleton), the realistic 16S generators, and the COT occupation-measure machinery.
\field{Deliverables} The layered route: level-set basins on the chart cover (an
identifiable component finder), the geometric corridor discriminator (overlap mass
$+$ elongated covariance), the local-covariance basin shape, the alignment of
per-subject basins onto shared components, and the transfer operator on the coarse
basin-and-corridor skeleton that supplies the dynamics and earns the
``quasi-equilibrium'' label. Spec: the COT bridge note.
\field{Gate} \begin{gate} On synthetic 16S with \emph{known} components: the
geometry-first level-set basins are more stable across scales and bootstraps than
the dictionary-factorization atoms; the transfer-operator metastable count agrees
with the cluster-tree count; the corridor discriminator removes spurious components.
Then, and only then, the HMP/152-subject experiment.\end{gate}
\field{Unblocks} The terminal application.

%==============================================================================
\section{Terminal --- the real 16S / omics application}\label{sec:terminal}
%==============================================================================

\field{Objective} Do the thing the program exists for: estimate conditional
expectations of binary and continuous features, and the community
occupancy/component structure, on real 16S and other omics data, with honest
uncertainty.
\field{Inputs} All prior phases.
\field{Deliverables} Real-data conditional-expectation surfaces with bands;
occupancy/component decompositions; reported under nested and grouped (subject-level)
cross-validation so the headline numbers are not selection-optimistic or
leakage-inflated.
\field{Gate} \begin{gate} Held-out predictive validity on real data and biological
interpretability of the components --- judged, as throughout, by prediction, not by
intrinsic-geometry scores.\end{gate}

%==============================================================================
\section{Where this diverges from the initial ordering, and why}\label{sec:diverge}
%==============================================================================

The initial list was: (1) LPS validation, (2) local-dimension regularization,
(3) PS-LPS, (4) covariance $+$ simplicial complex $+$ synthetic data, (5) occupancy.
The spine of that ordering is right and is preserved. The refinements are:

\begin{enumerate}[leftmargin=1.6em]
\item \textbf{Validation first --- agreed, emphatically}, and for a stronger reason
than ``it gates the rest'': Phase 0 \emph{emits} the $S$/df/LOO/bands toolkit that
the df-matched PS-LPS comparison, the dimension field's GCV, and the covariance
bands all reuse. It is the foundation that hands the later phases their tools.
\item \textbf{Data enabler pulled forward} (\S\ref{sec:data}). It was bundled near
the end of the list; it is in fact a dependency of items (2), (4), and (5), so it
must lead.
\item \textbf{PS-LPS and the dimension field run in parallel}, not in sequence.
They are independent extensions of the validated base; the list put one before the
other, but nothing connects them until they are combined in Phase 2b.
\item \textbf{The simplicial complex is the integration phase} (2b), not a sub-part
of the covariance work. It is where synchronization, dimension, and covariance are
unified over the nerve.
\item \textbf{The terminal real-data application is named explicitly} as the
endpoint, and \textbf{the audit-independence workflow wraps every phase} as the
governance layer.
\end{enumerate}

%==============================================================================
\section{Decision gates that can re-route the program}\label{sec:reroute}
%==============================================================================

The plan is adaptive, not a fixed waterfall: several gates can prune whole branches.

\begin{itemize}[leftmargin=1.4em]
\item \textbf{After E2 (data).} If the nominal$-$effective dimension gap in the real
VALENCIA strata is small, the within-face dimension field (1b) is deprioritized to
the cap-and-soft-face refinement, and 2a/3 proceed without it.
\item \textbf{After T9 (1b).} If the dimension field does not improve Truth-RMSE, it
is not promoted to production; a global chart dimension is kept.
\item \textbf{After 2b.} If higher-order ($2$-simplex) synchronization does not
measurably help, the program stays on the nerve $1$-skeleton.
\item \textbf{After Phase 3 on synthetic.} If the geometry-first basins are not more
stable than the existing factorization, the capstone falls back to the COT
program's current dictionary route rather than replacing it.
\end{itemize}

%==============================================================================
\section{Immediate next actions}\label{sec:next}
%==============================================================================

\begin{enumerate}[leftmargin=1.6em]
\item \textbf{Audit the LPS Tier-0 run} --- does the correctness battery pass, and
is the linear-smoother toolkit correct? (Queued; you will share the results.)
\item \textbf{Audit the LPS binary GM/FF run} --- the binary-mode behaviour, with the
fallback-telemetry and selection-metric fixes from the earlier results audit.
\item \textbf{In parallel}, confirm the toy stratified generators are ready, so
Phase 1b's boundary and scale tests have their known-truth inputs the moment Phase 0
clears its gate.
\end{enumerate}

Once these two audits are adjudicated, Phase 0 is gated and the parallel Phase-1
tracks (PS-LPS and the dimension field) begin together.

%==============================================================================
\begin{thebibliography}{9}\small
\bibitem{plan} \emph{LPS Correctness and Validation: A Precise Experimental Plan}
  (Phase 0 spec),
  \nolinkurl{~/current_projects/geosmooth/project_briefs/lps_experimental_plan_2026-06-09}.
\bibitem{audit} \emph{LPS/PS-LPS JASA-style audit memo} (PS-LPS linearity, df,
  $\lambda_{\mathrm{sync}}$ selection; the binary-run results audit),
  \nolinkurl{~/current_projects/geosmooth/project_briefs/lps_ps_lps_jasa_style_audit_2026-06-09}.
\bibitem{impl} \emph{LPS Local Dimension Regularization: Implementation Notes}
  (Phase 1b spec; the data enabler E1/E2/E3 and synthetic generators),
  \nolinkurl{~/current_projects/geosmooth/project_briefs/lps_local_dimension_regularization_implementation_notes_2026-06-10}.
\bibitem{chartaware} \emph{Chart-Aware Local Association Structure on the
  $L^\infty$ Atlas} (Phase 2a spec),
  \nolinkurl{~/current_projects/geosmooth/project_briefs/chart_aware_local_association_design_2026-06-09}.
\bibitem{nerve} \emph{The Local-Chart Simplicial Complex as an Organizing Substrate
  for LPS} (Phase 2b spec),
  \nolinkurl{~/current_projects/geosmooth/project_briefs/lps_chart_nerve_substrate_2026-06-10}.
\bibitem{bridge} \emph{Occupancy Supports as Basins on the Chart Cover} (Phase 3
  spec),
  \nolinkurl{~/current_projects/vaginal_community_trajectory_types/docs/occupancy_basins_on_chart_cover_bridge_2026-06-10}.
\bibitem{workflow} \emph{Two-Agent Research Project Workflow} (the governance
  wrapper / audit-independence rules),
  \nolinkurl{~/.codex/notes/workflows/two_agent_research_project_workflow}.
\end{thebibliography}

\end{document}
